% Source: https://artofproblemsolving.com/community/c3h1569746_sum_of_numbers_with_5_factors
% Stable identity: aops-topic-1569746-q01
If $n=p_1^{e_1}p_2^{e_2}\cdots p_r^{e_r}$, the number of divisors of $n$ is $(e_1+1)(e_2+1)\cdots(e_r+1)$. Since $5$ is prime, this product equals $5$ only when $r=1$ and $e_1=4$: the numbers with exactly $5$ divisors are the fourth powers of \emph{primes}, $p^4$.

The primes with $p^4<10000$ are $p=2,3,5,7$ ($11^4=14641$ is too big). Hence the sum is
\[
16+81+625+2401=\boxed{3123}.
\]
The sum $2^4+3^4+\cdots+9^4=15332$ also includes the fourth powers of $4,6,8,9$, which are not primes; for example $4^4=2^8$ has $9$ divisors and $6^4=2^4\cdot3^4$ has $25$.
